% 第二十版独立法英复校；原页与删改记录见 V20-D01--V20-D11。
% 起点：法文物理69页 Proposition 6.1；英文46页 Proposition 6.1。
\subsec{第6--8节的推广指示}
\source{法文手稿物理第69--70页；英文编辑稿第46页}

命题6.1中，应读作
\[
\dim f(X_i)\geq m+1,
\]
并在后文读作
\[
\dim f(Z)\leq m-1.
\]

7.1的一般论述对于任意$m$仍可照原样保留。7.2、7.3作必要的相应修改即可（请留意7.2中的记号$m$，它在这里容易造成混淆）。另一方面，在7.4的证明中，不再需要逐步推进，而可以直接取一个余维$m=n$的线性截面。\cmt{V20-D01}{以上是原作者对各条的编写指示，不是删去原假设后的完整命题。6.1的固有性、连通性等条件仍须保留；像集的维数依原条目的约定理解。$m=n$仅指7.4证明中截面余维的选择。法文的\emph{laïus}就是“一般论述”，英稿附加的“Fr or Latin”“speech”及页码69不是数学条目。旧注D17-M025、D17-M026的这些限制继续保留。}

在8.2中，把其中的超额维数条件改为
\[
\dim Y_\xi>\dim X-m,
\]
并把假设$\dim f(X_i)>0$替换为$\dim f(X_i)\geq m$。8.2的后续部分作类似修改。既然8.3给的是一个例子，就没有必要改动它，仍取$m=1$。\cmt{V20-D02}{法文第69页和英文第46页都把被替换的旧条件写成$\dim Y_\xi>\dim X$；但8.2原条目所列的是$\dim Y_\xi>\dim X-1$（英文第21页）。故正文以“其中的超额维数条件”指回旧项，不重抄这个遗漏。法文新条件末项有叠写痕迹，英稿排为$m$并保留“illegible, ask AG”；本次采用与余维$m$相应的阈值$\dim X-m$，但不把叠写字形宣称为完全清晰。这里的修改指示也不能消除8.2原假设自身尚有的读字问题；旧注D17-M027的这一保留继续有效。}

我把这个大概还须多费一点心思的练习留给你：找出对应于8.4、8.5、8.6（第30、31、32页）、适用于任意$m$的恰当陈述。其实，如果做这道练习使你厌烦，想来也并非非做不可。\cmt{V20-D03}{法文第69页手添页码30，页末两行可读为\emph{il n'est d'ailleurs sans doute pas indispensable de faire cet exercice, s'il te barbe de le faire}；正文保留其劝退语气。英稿在“你”之后另添“Dieudonné or Blass [Tr]”，这是英译者对受话者的猜测，并非法稿署名或作者原注。所引页码属于原稿内部指页，不是本中文本页码。}

我认为，第8节自8.6起的各项论述，基本上也能推广到任意$m$的线性截面。正式写出这些内容，无疑会是一项相当冗长而繁琐的练习。我得承认，我还不知道有哪一种应用在本质上依赖于对这一更一般情形的分析，因此我们并非一定要把这些讨论收入这部《基础》。另一方面，经验表明，在更一般的背景中写作，有时会迫使人把问题拆解得更好，从而更清楚地理解整套东西，而且通常并不会增加多少麻烦。此外，在这里这样一个真正的几何背景中，做一些形式表述方面的练习也并无坏处。当然，完全不能排除有朝一日我们会用到它，到时能找到现成的论述，会很高兴。不过，关于这件事的决定仍完全由你作出；我只把我们或许能就此给出的陈述概括一下。\cmt{V20-D04}{法文第70页清楚写\emph{à partir de 8.6}，并手添\emph{également}；故应为“自8.6起……也能”，而不是英稿的“except 8.6”。旧中文D17-M029据英稿译作“除8.6以外”，本次据法稿订正。这不是原作者前后自相矛盾。法稿页首另有两行被密集划销的旧文字，未能可靠恢复；不把它们拼成保留正文。“拆解”“整套东西”“通常并不会增加多少麻烦”分别保留\emph{dévisser}、\emph{le fourbis}、\emph{sans beaucoup plus de mal}的写作语气。}

\subsec{奇异零点与交维数的过滤}
\source{法文手稿物理第70页下半至第71页中部；英文编辑稿第46页下半至第47页上半}

再假设$f\colon X\to P$无分歧，且$X$在$S$上光滑，各分支的相对维数至少为$m$。在$X^{(m)}=V^0$中，区分出子预概形
\[
X^{(m)\mathrm{sing}}=V^1,
\]
它由$\phi^m$相对于$\operatorname{Grass}^{m}$的奇异零点组成。从几何上看，它也由这样的点对$(x,\xi)$组成：线性子簇$L_\xi$与$X$在$x$处的切空间（视为$P$的线性子簇）超额相交，也就是说，这两个空间不能张成整个$P$。\cmt{V20-D05}{“各分支的维数”在这里按相对于$S$的几何纤维理解，$n$须在相对维数恒定的部分分别取值。这里的奇异零点是截面的微分秩下降处，不是关联总空间的绝对奇异点集，也不能无条件等同于每个截面纤维自身的非光滑点集。以子式定义时，$V^1$是$\operatorname{Fitt}_{n-m}(\Omega_{X^{(m)}/\operatorname{Grass}^{m}})$定义的闭子预概形；其与基变换的相容性及“奇异零点”与超额维数光滑纤维的区别，详见同号核查。该说明补充旧注D17-M030，并未恢复第16段补篇的全部定义。}

与$m=1$时的情形相反，对于任意$m$，态射
\[
X^{(m)\mathrm{sing}}\longrightarrow X
\]
一般不光滑。原因是：经过$x$并与给定的线性子簇$T\ni x$超额相交的那些$L$所组成的簇，一般并不在底域上光滑；这个簇只是下述光滑子簇的闭包：它由那些使$T\cap L_\xi$的维数恰好比“正常”维数$n-m$大一的$\xi$组成，其中$n=\dim T$，$m=\operatorname{codim}L$。\cmt{V20-D06}{法文第70页末的手添语为\emph{passant par x et qui}，第71页首清楚接为\emph{est seulement l'adhérence de la sous-variété lisse}。因此，英稿的两处难辨标记及“only the loneve”残词均可据法稿恢复；旧注D17-M031不再以缺文保留这两处。闭包是前述光滑子簇的闭包，原句并未称闭包本身光滑；这里也不由底层集合的闭包关系推断任意概形结构的相等。英稿单列的“V [page 71 Fr]”在法稿中对应公式后的手写勾记，不作为新对象或条目引用。}

如果没有弄错，第16段补篇中引入的集合$V''$（它包含在相对超奇异点集中），无非就是由满足
\[
\dim(T_x\cap L_\xi)\geq n-m+2
\]
的点对$(x,L_\xi)$组成的集合。因此，$V^1\setminus V''$在$X$上光滑；如果没有弄错，它甚至恰好是$V^1$在$X$上光滑的那些点所组成的集合。\cmt{V20-D07}{法文第71页两处均为打字上标$1$的$V^1$，而且第一处明确是\emph{lisse sur X}。英文第47页改排为$V'$，并把第一处基底误作$S$；在当前假设下从$X$上的光滑性可推出$S$上的光滑性，但不能因此说英译逐字忠实。正文据法稿恢复$V^1$及$X$，英稿的$V'$仅作为同一对象的别写记录。Dot第119批另称英稿写成$V^0$，与实际英文页面不符，本版不采纳这一归责。两次“如果没有弄错”继续保留，特别是对补篇所定义$V''$的认同仍属原作者的待核陈述。}

（要验证这一点，需要按照交的维数研究格拉斯曼概形的过滤，其中$L$变动而$T$固定。如果没有弄错，对于从首个非平凡退化项开始的闭过滤，下一阶恰好由前一阶的非光滑点组成\ednote{英文底本脚注61：“询问A.G.”；正文在\emph{notch}之后还保留“Fr”“stratum???”。法文原词\emph{cran}在此是过滤中的一阶。}；这里须不仅在集合意义下，而且还借助第16段补篇第16页的引理\ednote{英文底本脚注62：“A.G.，请帮助定位这一引用。”}，在概形意义下定义这一过滤。因此，这项研究将构成专门讨论格拉斯曼簇的一个“几何性”段落中的一小节。）\cmt{V20-D08}{原稿未注明这句话从哪一阶开始。若把全格拉斯曼概形$V^{(0)}$也算作前一阶，就不成立：它本来光滑，而$V^{(1)}$可能非空。正文加入“从首个非平凡退化项开始”这一限制。对于自然子式结构，在$1\leq k\leq\min(m,r-1-n)$内可证$V^{(k)}\to X$的非光滑点支集恰是$V^{(k+1)}$的支集；核查报告给出一般矩阵坐标证明，未凭有限个例子作推广。原文说的是\emph{points non lisses}，不能加强成某个典范奇异子概形与下一退化概形的结构相等；报告另以$2\times3$矩阵说明两种概形结构可以不同。所引补篇引理本次仍未定位。}

\subsec{退化层的维数及其像}
\source{法文手稿物理第71页中下部，止于“on peut oublier V''”；英文编辑稿第47页中部}

若同样将$V^{(k)}$定义为$X^{(m)}$中对应于条件
\[
\dim(T_x\cap L)\geq n-m+k
\]
的子概形，则在一个固定的非空、纯维$n$的几何纤维上，设
\[
N=r-1,\qquad 0\leq m\leq n\leq N,
\]
当整数$k$满足
\[
0\leq k\leq\min(m,N-n)
\]
时，该退化层非空，而且维数计算给出
\[
\dim\operatorname{Grass}^{m}(P)-\dim V^{(k)}
=(k-1)(n-m)+k^2.
\]
（注意：更一般地，这来自上面所说的格拉斯曼簇过滤中各“胞腔”的维数计算。）\cmt{V20-D09}{法文第71页原写“至少在$k\leq m$、$k\leq r-m$这些合理限制下，除非计算有误”；英稿照录。这两个界不足以保证非空。取$X=P=\mathbf P^2_K$、$f=\id$、$r=3$、$n=2$、$m=k=1$，两项原界均成立，但每条过$x$的直线与$T_x=P$的交维数只有1，达不到$n-m+k=2$，故$V^{(1)}=\varnothing$。正文保留第十九版D17-M034及D17-M-E02已经得到的正确非空范围；本次重新推导而非仅复制旧结论。固定$x$时相应闭层的余维为$k(n-m+k)$，在纯维$n$的几何纤维上再加$n$，即得显示公式。各分支维数不同时须分别取$n$。法稿正式公式之前另有被密集划销的旧维数式，可辨开头为$\dim V^{(k)}$，右端不能可靠恢复；删稿身份及图像位置记于V20-D11，不据最后公式反填。}

若采用上面对$V''$的集合识别，则在上述几何纤维为有限型的情形下，取$k=2$，非空时得到的维数差为$n-m+4\geq4$，因此$V''$在$\operatorname{Grass}^{m}$中的像的闭包余维至少为4；若二阶退化层为空，则无须排除它的像。于是，如果我们只关心去掉$\operatorname{Grass}^{m}$中余维至少为2的子集之后的情形，就可以不考虑$V''$。\cmt{V20-D10}{法文在这里的“像”“余维至少4”以及手添的参数空间\emph{de Grass}$^m$都可辨；英稿在Grass之后的难辨标记不再沿入正文。维数不等式需要像闭包的维数不超过源维数，本次明确采用域上有限型的条件，并写“像的闭包”，不另假设像已经闭。若$X$固有，则相应投影固有，像确实闭；一般情形不能省去这个区分。对$V''$的结论仍以作者前面的集合认同为条件，并不因此验证其超奇异概形结构。旧注D17-M035的代入计算保留，但本次已有一般的退化层维数证明及像闭包论证。}
% 精确止点：法文物理71页“on peut oublier V''”。
% 下句“D'autre part, dans X_{Grass^m} - V''”由后续模块承接，本文件不含。
